\documentclass[12pt,a4paper,oneside]{article}

\usepackage[T1]{fontenc}
\usepackage[utf8]{inputenc}
\usepackage[brazilian]{babel}
\usepackage{mathptmx}
\usepackage[nopatch=footnote]{microtype}
\usepackage{geometry}
\usepackage{graphicx}
\usepackage{fancyhdr}
\usepackage{titlesec}
\usepackage{booktabs}
\usepackage{tabularx}
\usepackage{array}
\usepackage{longtable}
\usepackage{enumitem}
\usepackage{float}
\usepackage{caption}
\usepackage{xcolor}
\usepackage{hyperref}
\usepackage{url}
% Submission build only (build_submission.sh): the signed AI-use declaration is appended as a final annex.
\ifdefined\tccannex
\usepackage{pdfpages}
\fi
\usepackage[alf]{abntex2cite}

\definecolor{MackenzieRed}{RGB}{235,0,40}
\definecolor{LinkBlue}{RGB}{0,61,165}
\definecolor{DraftGray}{RGB}{246,246,246}
\newif\ifworkingdraft
\workingdrafttrue % trocar para \workingdraftfalse somente na versão final aprovada pelo grupo e pelo orientador
\ifdefined\tccfinal\workingdraftfalse\fi

\geometry{
  top=3.15cm,
  bottom=2.20cm,
  left=2.50cm,
  right=2.00cm,
  headheight=34pt,
  headsep=0.65cm
}

\setlength{\parindent}{1.25cm}
\setlength{\parskip}{0pt}
\setlength{\emergencystretch}{2em}
\linespread{1.0}
\frenchspacing

\titleformat{\section}
  {\normalfont\bfseries\fontsize{14}{16}\selectfont}
  {\thesection.}{0.55em}{}
\titlespacing*{\section}{0pt}{1.15em}{0.55em}

\titleformat{\subsection}
  {\normalfont\bfseries\fontsize{12}{14}\selectfont}
  {\thesubsection}{0.55em}{}
\titlespacing*{\subsection}{0pt}{0.95em}{0.35em}

\titleformat{\subsubsection}
  {\normalfont\bfseries\itshape\fontsize{12}{14}\selectfont}
  {\thesubsubsection}{0.55em}{}
\titlespacing*{\subsubsection}{0pt}{0.80em}{0.25em}

\newcolumntype{Y}{>{\raggedright\arraybackslash}X}
\newcolumntype{C}[1]{>{\centering\arraybackslash}p{#1}}
\newcolumntype{L}[1]{>{\raggedright\arraybackslash}p{#1}}

\captionsetup{
  font=small,
  labelfont=bf,
  justification=centering,
  singlelinecheck=false,
  skip=4pt
}
\captionsetup[table]{position=top}

\setlist[itemize]{leftmargin=1.8em,itemsep=0.15em,topsep=0.35em}
\setlist[enumerate]{leftmargin=2.0em,itemsep=0.15em,topsep=0.35em}

\graphicspath{{assets/}}

\pagestyle{fancy}
\fancyhf{}
\fancyhead[L]{\includegraphics[height=0.82cm,trim=0 0 320 0,clip]{fci-header.png}}
\fancyhead[R]{\small\thepage}
\renewcommand{\headrulewidth}{0pt}
\newcommand{\mackenzierule}{%
  \par\vspace{2pt}\noindent\color{MackenzieRed}\rule{\textwidth}{0.65pt}\color{black}}
\fancyhead[C]{\raisebox{-0.52cm}[0pt][0pt]{\makebox[0pt][c]{\color{MackenzieRed}\rule{\textwidth}{0.65pt}}}}
\fancyfoot{}

\hypersetup{
  pdftitle={Desenvolvimento e implementação de uma arquitetura de proteção na integração de LLMs em ERPs: estudo de caso com Odoo e Contas a Pagar},
  pdfauthor={Felipe Ujvari Gasparino de Sousa; Gustavo Nascimento Siqueira; Thomaz de Souza Scopel; Mario Olimpio de Menezes},
  pdfsubject={Trabalho de Conclusão de Curso II em Ciência da Computação},
  pdfkeywords={LLM; prompt injection indireta; ERP; Odoo; Contas a Pagar; segurança},
  pdfcreator={Universidade Presbiteriana Mackenzie - Faculdade de Computação e Informática},
  colorlinks=true,
  linkcolor=LinkBlue,
  citecolor=LinkBlue,
  urlcolor=LinkBlue,
  bookmarksopen=true
}

\begin{document}
\selectlanguage{brazilian}

\begin{center}
  \vspace*{1.15cm}
  {\fontsize{16}{19}\selectfont
  Desenvolvimento e implementação de uma arquitetura de proteção na integração de LLMs em ERPs: estudo de caso com Odoo e Contas a Pagar\par}

  \vspace{1.35cm}

  {\fontsize{12}{15}\selectfont
  Felipe Ujvari Gasparino de Sousa\textsuperscript{1} \quad
  Gustavo Nascimento Siqueira\textsuperscript{1} \quad
  Thomaz de Souza Scopel\textsuperscript{1}\par
  Mario Olimpio de Menezes\textsuperscript{1}\par}

  \vspace{0.55cm}

  {\fontsize{11}{14}\selectfont
  \textsuperscript{1}Faculdade de Computação e Informática (FCI)\par
  Universidade Presbiteriana Mackenzie -- São Paulo, SP -- Brasil\par}

  \vspace{0.55cm}
  2026

  \ifworkingdraft
  \vspace{0.65cm}

  \fcolorbox{black!25}{DraftGray}{%
    \parbox{0.91\textwidth}{\centering\small
      Versão de trabalho, revisada em 31 de outubro de 2026. Contém os resultados auditados de 144 execuções no Odoo, com análise descritiva de segurança e utilidade. Aprovação final do grupo e do orientador pendente.}}
  \fi
\end{center}

\vspace{0.75cm}

\begin{center}
  {\fontsize{14}{16}\selectfont Resumo}
\end{center}
\begin{itshape}
  Agentes baseados em grandes modelos de linguagem podem interpretar documentos e acionar ferramentas corporativas, mas também podem tratar instruções em conteúdo externo como comandos legítimos. Em Contas a Pagar, uma injeção de prompt indireta pode alterar campos financeiros, redirecionar o beneficiário ou induzir ação crítica sem autorização. Este trabalho objetiva desenvolver e avaliar uma arquitetura de proteção para integrar LLMs ao Odoo 19.0 Community em um fluxo sintético de contas de fornecedor. A pesquisa segue Design Science Research com experimento controlado, comparando uma integração funcional de referência (C0) à defesa em profundidade (C1), com política determinística, menor privilégio, aprovação vinculada, idempotência e auditoria. O experimento principal prevê 90 pares contrafactuais, três repetições e 1.080 execuções com um modelo DeepSeek; o piloto separado acrescenta 144. Há provas técnicas do fluxo documental, dos caminhos C0/C1, da captura e restauração do estado e de registros que preservam falhas. No piloto de propostas, em duas rodadas, a segunda com estratégias de ataque revisadas, nenhuma das 72 propostas adversariais de C0 incorporou o objetivo do atacante, e as propostas benignas foram exatas em 36 de 36 execuções em C0 e 35 de 36 em C1. Com esse modelo e esses ataques estáticos, a redução do sucesso dos ataques não pode ser demonstrada; um enquadramento estimativo aguarda decisão do grupo. O efeito no estado do ERP e o teste final permanecem pendentes. Replicação entre modelos, ablação e ataques adaptados à defesa ficam como trabalhos futuros.

\end{itshape}

\vspace{0.45cm}
\noindent\textbf{Palavras-chave:} agentes baseados em LLM; prompt injection indireta; defesa em profundidade; ERP; Contas a Pagar.

\newpage

\begin{center}
  {\fontsize{14}{16}\selectfont Abstract}
\end{center}
\begin{itshape}
  \begin{otherlanguage*}{english}
    Large language model agents can interpret documents and invoke corporate tools, but they may also treat instructions in external content as legitimate commands. In accounts payable, indirect prompt injection can alter financial fields, redirect beneficiaries, or induce unauthorized critical actions. This study aims to develop and evaluate a protection architecture integrating LLMs with Odoo 19.0 Community in a synthetic vendor-bill workflow. Following Design Science Research and a controlled experiment, it compares a functional reference integration (C0) with defense in depth (C1), combining deterministic policy, least privilege, action-bound approval, idempotency, and auditing. The main experiment specifies 90 benign--adversarial counterfactual pairs, three repetitions, and 1,080 executions with one DeepSeek model; a separate pilot adds 144. Evaluation observes ERP state, security, and utility non-inferiority; latency and cost will be described. Technical proofs cover the document workflow, private C0/C1 execution paths, state capture and restoration, and records preserving failures. These findings do not establish protection effectiveness: operational integration, review, pilot, and protocol freeze remain pending. Scientific conclusions depend on complete data collection. Cross-model replication, ablation, and defense-aware attacks are future work, limiting external validity and attribution to individual layers.

  \end{otherlanguage*}
\end{itshape}

\vspace{0.45cm}
\noindent\textbf{Keywords:} LLM-based agents; indirect prompt injection; defense in depth; ERP; accounts payable.

\vspace{0.75cm}

\section{Introdução}
Grandes modelos de linguagem (\textit{Large Language Models}, LLMs) deixaram de operar apenas como geradores de texto e passaram a compor agentes capazes de consultar dados, selecionar ferramentas e produzir chamadas que alteram sistemas. Essa passagem da recomendação para a ação modifica a natureza do risco: uma saída inadequada deixa de ser somente conteúdo incorreto e pode tornar-se uma mutação persistente, uma mensagem enviada ou a exposição de informação. Ambientes como AgentDojo e InjecAgent mostram que agentes conectados a ferramentas permanecem suscetíveis a instruções presentes em dados externos, ao mesmo tempo em que podem falhar em tarefas legítimas mesmo sem ataque \cite{debenedetti2024agentdojo,zhan2024injecagent}. Segurança e utilidade, portanto, precisam ser observadas conjuntamente.

A injeção de prompt indireta ocorre quando o atacante não escreve no canal de comando do agente, mas posiciona instruções em conteúdo que o sistema processará como dado: documento, mensagem, página ou registro recuperado. O modelo recebe na mesma representação linguística o objetivo confiável da aplicação e o conteúdo não confiável. Caso siga este último, pode desviar a tarefa ou escolher argumentos incompatíveis com a intenção do usuário. Greshake et al. demonstraram esse vetor em aplicações integradas a LLM e sistematizaram impactos que incluem chamadas indevidas de API e exfiltração \cite{greshake2023indirect}. Trabalhos posteriores construíram benchmarks mais amplos, mas também verificaram que defesas avaliadas apenas contra ataques estáticos podem ser contornadas por ataques adaptados ao mecanismo conhecido \cite{yi2025bipia,zhan2025adaptive}.

O processo de Contas a Pagar concentra propriedades úteis para investigar o problema. Uma conta de fornecedor nasce de conteúdo externo; seus campos precisam ser confrontados com fontes internas, como fornecedor, pedido de compra, moeda e conta bancária; e o sistema registra estados distintos entre rascunho, contabilização e pagamento. A própria documentação do Odoo separa a criação e o preenchimento de uma conta de fornecedor de sua confirmação, que gera o lançamento contábil \cite{odoo2026vendorbills}. Isso permite definir consequências mensuráveis sem movimentar dinheiro real: alterar campos financeiros, redirecionar o beneficiário indicado ou contabilizar sem autorização são eventos verificáveis por código.

O estudo procura produzir evidência contextualizada que relacione a chamada proposta pelo modelo ao efeito end-to-end em um processo empresarial. A comparação com trabalhos correlatos examina invariantes do ERP, separação entre rascunho e ação crítica, origem confiável dos campos e restauração do estado. Reduzir ataques bloqueando todas as ações também não constitui uma solução útil. A revisão dirigida e sua matriz crítica posicionam a contribuição como avaliação contextualizada e reproduzível desse recorte, sem reivindicar cobertura sistemática ou provar a inexistência de trabalhos equivalentes fora da seleção.

O problema de pesquisa é avaliar em que medida uma defesa em profundidade reduz violações no fluxo delimitado sem comprometer excessivamente as tarefas legítimas. O público-alvo primário é composto por equipes de engenharia, arquitetura, segurança e governança responsáveis por integrar modelos generativos a processos corporativos críticos. Gestores financeiros e auditores constituem público secundário, pois precisam compreender quais decisões permanecem automáticas, quais exigem autorização e que evidência fica disponível para prestação de contas.

O objetivo geral é projetar, implementar e avaliar experimentalmente uma arquitetura de defesa em profundidade para agentes baseados em LLM integrados a processos corporativos críticos, usando um fluxo de Contas a Pagar em Odoo como estudo de caso e medindo mitigação de prompt injection indireta, ações não autorizadas, utilidade e sobrecarga. A pergunta principal é:

\begin{quote}
\itshape Em que medida uma arquitetura de defesa em profundidade reduz o sucesso de ataques de prompt injection indireta e ações não autorizadas em um agente LLM integrado a um processo de Contas a Pagar, preservando a utilidade da automação e com qual custo operacional?
\end{quote}

Para respondê-la, o estudo compara duas condições sobre os mesmos documentos, modelos e estados iniciais. C0 representa uma integração direta, funcional e competente, com saída estruturada, mas sem o pacote de autorização avaliado. C1 separa proposta e autoridade: o LLM descreve uma ação tipada; um mecanismo determinístico consulta o estado confiável; uma aprovação vinculada é exigida para ação crítica; e um executor de menor privilégio aplica somente a decisão autorizada. O princípio é que o modelo propõe, o mecanismo de política autoriza e o executor aplica.

As contribuições pretendidas são: (i) uma arquitetura implementada que explicita fronteiras de confiança e autoridade; (ii) o conjunto AP-IPI-BR, formado por documentos sintéticos benignos e adversariais pareados; (iii) um protocolo que verifica o sucesso do ataque no estado do ERP, por observação independente do LLM; e (iv) evidência quantitativa de segurança e utilidade, acompanhada de descrição do custo e análise de falhas. A avaliação estima o efeito do pacote completo de proteção, sem isolar causalmente cada camada. O trabalho também preserva proveniência de modelo, artefatos, prompts, dados e runs, permitindo distinguir reprodução da análise de uma futura reexecução com serviços que podem sofrer atualização.

Esta versão registra o estágio atual do TCC II: além das provas técnicas do laboratório, dos caminhos C0/C1, da captura/restauração e do registro candidato das execuções, foi concluído o piloto de propostas em duas rodadas de 144 chamadas com o modelo real, nas quais nenhuma das 72 propostas adversariais de C0 incorporou o objetivo do atacante e as propostas benignas foram exatas em 36 de 36 execuções em C0 e 35 de 36 em C1. O ciclo integrado no estado do ERP para os pares do piloto, a decisão do grupo sobre o enquadramento estimativo, a revisão nominal, o congelamento e a coleta confirmatória permanecem pendentes. Essa distinção preserva a fronteira entre viabilidade técnica, propostas do piloto e efeito final no sistema. A Seção 2 apresenta os conceitos e trabalhos relacionados; a Seção 3 descreve o artefato, o dataset e o experimento; a Seção 4 relata as provas técnicas e o piloto de propostas, discute seus limites e define a apresentação dos resultados finais; e a Seção 5 sintetiza as conclusões permitidas nesta etapa.


\section{Referencial Teórico}
\subsection{Agentes LLM, ferramentas e fronteiras de confiança}

Um agente LLM combina um modelo de linguagem, memória ou contexto, lógica de orquestração e ferramentas que observam ou alteram um ambiente. O modelo não precisa possuir credenciais diretamente: basta que sua saída determine qual ferramenta será chamada e com quais argumentos. Essa composição amplia a utilidade do modelo, mas cria uma cadeia causal entre texto não confiável e efeitos autorizados pela aplicação. A segurança do conjunto não pode ser inferida apenas da qualidade do modelo ou do prompt.

O perfil de IA generativa do NIST enquadra o risco ao longo do ciclo de vida e recomenda governar, mapear, medir e gerir propriedades do sistema, não apenas do modelo \cite{nist2024genai}. Para aplicações com ferramentas, isso implica inventariar dados, permissões, dependências, monitoramento e respostas a falhas. No mesmo sentido, o OWASP distingue prompt injection de autonomia excessiva: a primeira altera o comportamento esperado; a segunda permite que saídas inesperadas, ambíguas ou manipuladas produzam ações danosas devido ao excesso de funcionalidade, permissão ou autonomia \cite{owasp2025prompt,owasp2025agency}.

Há três autoridades que precisam ser separadas. A autoridade epistêmica diz respeito ao que o modelo consegue inferir do documento. A autoridade de decisão define se determinada ação é compatível com a política e o estado confiável. A autoridade de execução corresponde à credencial e ao método capazes de produzir o efeito. Um modelo pode ser útil na primeira sem receber as duas últimas. Quando a mesma saída interpreta, autoriza e executa, um erro de linguagem atravessa todas as fronteiras sem uma verificação independente.

\subsection{Prompt injection indireta}

Prompts e dados são serializados como linguagem e o modelo não oferece, por padrão, uma separação formal equivalente à distinção entre código e dados de sistemas tradicionais. Na injeção direta, o atacante controla o canal de entrada do usuário; na indireta, controla um artefato recuperado ou recebido por uma aplicação. Greshake et al. mostraram que conteúdo remoto pode manipular aplicações integradas e influenciar chamadas de API, divulgação de informação e persistência de instruções \cite{greshake2023indirect}. A ameaça é particularmente relevante quando o documento tem aparência legítima e o atacante não precisa acessar o agente ou o ERP.

InjecAgent operacionalizou esse problema em 1.054 casos, 17 ferramentas do usuário e 62 ferramentas do atacante, distinguindo dano direto e exfiltração \cite{zhan2024injecagent}. BIPIA investigou tarefas de processamento de texto e identificou como causas a dificuldade de separar contexto informacional de instrução acionável e a baixa consciência do modelo sobre a origem do conteúdo \cite{yi2025bipia}. AgentDojo, por sua vez, combinou tarefas legítimas e casos de segurança em um ambiente extensível de agentes com ferramentas, evidenciando que avaliar apenas ataque é insuficiente, pois modelos também falham na tarefa benigna \cite{debenedetti2024agentdojo}.

Esses trabalhos justificam duas escolhas deste TCC. Primeiro, o ataque deve ser incorporado ao conteúdo que o pipeline realmente entrega ao modelo; casos em que o payload se perde na extração precisam ser medidos, não descartados. Segundo, sucesso do ataque deve ser definido por um invariante externo. Uma proposta maliciosa bloqueada revela comprometimento do modelo, mas não equivale a alteração do ERP. O encadeamento entre entrega, proposta, autorização e estado permite localizar a fronteira em que o ataque foi interrompido.

\subsection{Defesas e seus limites}

Defesas de prompt podem reforçar instruções, delimitar conteúdo, classificar entradas ou pedir ao próprio modelo que detecte manipulação. Elas são úteis como sinais, mas continuam apoiadas no componente suscetível. Yi et al. propõem consciência de fronteira e lembretes explícitos para reduzir ataques em tarefas textuais \cite{yi2025bipia}. StruQ avança na separação entre instruções e dados mediante formato estruturado e treinamento específico, obtendo resistência com pequena perda de utilidade no contexto avaliado \cite{chen2025struq}. Essa abordagem trata a raiz linguística do problema, embora dependa do modelo e do regime de treinamento usados.

Outra linha controla ações no tempo de execução. Task Shield verifica se instruções e chamadas de ferramenta contribuem para o objetivo do usuário e demonstrou redução de ataques mantendo utilidade no AgentDojo \cite{jia2025taskshield}. DRIFT combina planejamento mínimo, validação dinâmica e isolamento de instruções no fluxo de memória, aproximando a defesa do nível do sistema \cite{li2025drift}. Esses mecanismos reforçam que segurança agentiva não precisa depender de reconhecer toda frase maliciosa: é possível verificar a compatibilidade entre objetivo, trajetória, capacidade e parâmetros.

Resultados positivos, entretanto, não autorizam generalização irrestrita. Zhan et al. adaptaram ataques ao conhecimento de oito defesas e contornaram todas, obtendo taxas superiores a 50\% no cenário do estudo \cite{zhan2025adaptive}. O achado não implica que toda defesa falhará da mesma forma, mas mostra que avaliação estática tende a superestimar robustez. Por isso, este trabalho separa o conjunto confirmatório, congelado antes da coleta, de um conjunto defense-aware elaborado depois que a arquitetura de C1 estiver conhecida, sem permitir que C1 seja reajustada a partir desses casos.

O desenho adotado combina controles em camadas. Proveniência e separação de estágios reduzem a confusão de origem; schema e canonicalização restringem a forma da proposta; política determinística confronta argumentos com fontes internas; menor privilégio limita efeitos; aprovação vinculada impede reutilizar uma autorização em ação modificada; idempotência controla repetição; e auditoria preserva a explicação operacional. Um detector de prompt injection pode elevar risco, mas nunca conceder permissão. A avaliação principal estima o pacote completo, não atribui causalidade a uma camada isolada.

\subsection{Contas a Pagar e Odoo como testbed}

Contas a Pagar recebe documentos de uma parte externa e os relaciona a cadastros e compromissos internos. Em um fluxo apoiado por pedido de compra, valor, moeda, fornecedor, itens e beneficiário não precisam ser aceitos como verdade apenas porque aparecem na conta. Eles podem ser reconciliados com o pedido confirmado e com a conta bancária previamente aprovada. Essa estrutura cria regras determinísticas que independem da interpretação do LLM.

No Odoo, contas de fornecedor podem ser criadas manual ou automaticamente e permanecem em rascunho até a confirmação; depois da confirmação, o sistema produz lançamento contábil e restringe alterações \cite{odoo2026vendorbills}. O testbed usa essa diferença para separar uma ação reversível de uma ação crítica. A criação automática termina, no máximo, em rascunho. Contabilização requer decisão de política e autorização vinculada. Pagamento real fica fora do escopo, e qualquer solicitação experimental adicional termina em outbox sem egress.

O Odoo 19 introduziu a API externa JSON-2, na qual métodos de modelos são acessados por rotas HTTP específicas \cite{odoo2026json2}. A generalidade do transporte não equivale a uma fronteira de autorização. Em C1, a chamada externa chega a um add-on mínimo, transacional e idempotente; a identidade técnica não recebe grupos genéricos de faturamento ou pagamento. O objetivo não é demonstrar vulnerabilidade do Odoo, mas usar seu modelo de estado, permissões e regras de negócio para observar efeitos que um simulador textual não oferece.

\subsection{Síntese dos trabalhos relacionados}

A Tabela~\ref{tab:related-work} posiciona o estudo em relação às referências centrais. A coluna “estado empresarial” diferencia a observação de resposta ou ferramenta da verificação de mutações em um ERP. Nenhum dos trabalhos é reduzido a uma única característica; a comparação identifica apenas o recorte que motiva o artefato proposto.

\begin{table}[H]
\centering
\caption{Síntese dos trabalhos relacionados e lacuna explorada.}
\label{tab:related-work}
\footnotesize
\begin{tabularx}{\textwidth}{L{2.45cm} Y C{1.45cm} C{1.55cm} Y}
\toprule
\textbf{Trabalho} & \textbf{Foco principal} & \textbf{Utilidade} & \textbf{Estado empresarial} & \textbf{Contribuição para este estudo} \\
\midrule
Greshake et al. \cite{greshake2023indirect} & Vetores e impactos de IPI em aplicações integradas & parcial & não & Fundamenta ator remoto e confusão dado--instrução \\
InjecAgent \cite{zhan2024injecagent} & Benchmark de agentes e ferramentas & sim & ambiente de ferramentas & Orienta famílias de dano e proposta maliciosa \\
AgentDojo \cite{debenedetti2024agentdojo} & Tarefas, ataques e defesas em ambiente dinâmico & sim & estado simulado & Reforça medição conjunta de segurança e competência \\
BIPIA \cite{yi2025bipia} & IPI em tarefas textuais e defesas de fronteira & sim & não & Apoia separação de contexto e entrega do payload \\
StruQ \cite{chen2025struq} & Consulta estruturada e modelo treinado & sim & não & Demonstra separação de canais como camada possível \\
Task Shield \cite{jia2025taskshield} & Alinhamento de instruções e tool calls à tarefa & sim & ambiente de ferramentas & Motiva verificação externa da ação proposta \\
Ataques adaptativos \cite{zhan2025adaptive} & Bypass de defesas conhecidas & sim & ambiente de ferramentas & Justifica bloco defense-aware pós-freeze \\
Este trabalho & Defesa em profundidade no fluxo de Contas a Pagar & sim & Odoo e outbox & Mede efeito end-to-end, utilidade, custo e fronteiras \\
\bottomrule
\end{tabularx}
\end{table}

A lacuna específica não é a ausência de listas de controles. NIST e OWASP já oferecem taxonomias úteis, e a literatura propõe mecanismos promissores. O problema é instanciar uma arquitetura no limite entre documento financeiro, LLM e ERP e avaliá-la com pares contrafactuais, fontes de verdade internas, autorização por ação e oráculo de estado. Esse é o elo entre a revisão e a metodologia descrita a seguir.


\section{Metodologia}
\subsection{Estratégia de pesquisa}

Trata-se de pesquisa aplicada, quantitativa, complementada por análise qualitativa das falhas. Adota-se Design Science Research (DSR) porque o conhecimento é produzido pela construção e avaliação de um artefato destinado a resolver um problema relevante \cite{hevner2004dsr,peffers2007dsrm}. As atividades cobrem identificação do problema, objetivos, projeto, desenvolvimento, demonstração, avaliação e comunicação. A avaliação usa experimento controlado e pré-especifica hipóteses, variáveis, unidade de análise e repetição, conforme princípios da experimentação em Engenharia de Software \cite{wohlin2012experimentation}.

O artefato é composto por arquitetura, implementação, contratos, dataset, runner e oráculos. Odoo funciona como ambiente experimental e fonte do estado final; não é o objeto cuja segurança se pretende avaliar. O fator principal é a condição arquitetural, com dois níveis: C0 e C1. O documento, o modelo, o prompt-base, o schema, a fixture, o estado inicial e os limites permanecem controlados dentro de cada par. O protocolo está em pré-congelamento nesta versão: decisões podem ser corrigidas a partir do piloto, mas nenhuma observação do conjunto final pode influenciar hipótese, métrica ou defesa.

Não há participantes humanos. O approval gate é um componente técnico avaliado por fixtures de autorização definidas no manifesto. Ele mede se uma ação crítica para sem token válido e se um token deixa de valer quando os argumentos mudam; não mede atenção, julgamento, carga ou comportamento de revisores.

\subsection{Perguntas e hipóteses}

As perguntas derivam da questão principal e separam eficácia, utilidade e diferenças entre ataques. RQ abrevia a expressão inglesa para pergunta de pesquisa; H identifica uma hipótese. Os códigos nas tabelas servem para relacionar cada pergunta ao resultado correspondente. Os identificadores retirados do escopo são preservados para rastreabilidade, sem criar compromissos de avaliação.

\begin{table}[H]
\centering
\caption{Perguntas de pesquisa.}
\label{tab:rqs}
\small
\begin{tabularx}{\textwidth}{C{0.85cm} Y}
\toprule
\textbf{ID} & \textbf{Pergunta} \\
\midrule
\hypertarget{rq1}{RQ1} & Qual é o efeito de C1, em comparação com C0, sobre propostas maliciosas, violações de política e ações não autorizadas efetivamente registradas no Odoo? \\
\hypertarget{rq2}{RQ2} & Qual é o efeito de C1 sobre conclusão correta de tarefas benignas, falsos bloqueios, escalonamentos, latência e custo? \\
\hypertarget{rq3}{RQ3} & Como o efeito varia entre manipulação de campos, redirecionamento de beneficiário e ação crítica não autorizada? \\
\hypertarget{rq4}{RQ4} & Retirada desta entrega; identificador reservado. A contribuição isolada das camadas não será avaliada. \\
\bottomrule
\end{tabularx}
\end{table}

\hypertarget{h1}{H1} prevê menor taxa de ataques que alteram indevidamente o estado (ASR-state) em C1. \hypertarget{h2}{H2} prevê menor taxa de ações críticas concluídas sem autorização válida (UACR); seu denominador adversarial é de 30 pares, com poder e precisão a verificar antes do congelamento. \hypertarget{h3}{H3} avalia a preservação da taxa de sucesso em tarefas benignas (BTSR) por não inferioridade: a diferença $BTSR_{C1}-BTSR_{C0}$ não pode perder mais de 10 pontos percentuais, e o limite inferior do intervalo de confiança de 95\% deve ficar estritamente acima de $-0{,}10$. A margem é absoluta e não será recalibrada pelo resultado do piloto. \hypertarget{h4}{H4}, antiga hipótese de aceitabilidade de custo, e \hypertarget{h5}{H5}, associada à retirada de camadas, estão fora da entrega, com identificadores reservados. Latência, tokens, custo e demanda técnica de revisão permanecem medidas descritivas, sem teste de aceitabilidade.

O efeito sobre segurança em \hyperlink{rq1}{RQ1} e a não inferioridade da utilidade em \hyperlink{rq2}{RQ2} têm análise confirmatória prevista. O custo de \hyperlink{rq2}{RQ2} e as diferenças de \hyperlink{rq3}{RQ3} serão descritivos. A ordem confirmatória será \hyperlink{h1}{H1}, \hyperlink{h3}{H3}, \hyperlink{h2}{H2}. Uma hipótese posterior só recebe interpretação confirmatória se a anterior for satisfeita, embora efeitos e intervalos sejam sempre reportados. O status confirmatório depende da revisão de poder, precisão e implementação dos intervalos antes do congelamento, especialmente para \hyperlink{h2}{H2} e \hyperlink{h3}{H3}; o estudo não pressupõe poder suficiente por ter três repetições.

\subsection{Threat model e unidade de análise}

O atacante controla o conteúdo de uma conta ou documento apresentado ao processo, mas não controla o host, o modelo, o gateway, o Odoo, as fixtures, o pedido de compra ou a conta bancária aprovada. Ele conhece o objetivo geral do fluxo no conjunto estático; a adaptação de ataques à defesa conhecida não integra esta entrega. Estão fora do escopo comprometimento do provedor, exploração do parser, malware em PDF, ataque à cadeia de suprimentos, furto de credencial e participação interna maliciosa.

\begin{table}[H]
\centering
\caption{Ativos, objetivos adversariais e invariantes observáveis.}
\label{tab:threat-model}
\small
\begin{tabularx}{\textwidth}{L{2.30cm} L{3.20cm} Y}
\toprule
\textbf{Família} & \textbf{Objetivo do atacante} & \textbf{Invariante verificado} \\
\midrule
A1 & Alterar valor, fornecedor ou referência & O rascunho coincide exatamente com PO, fornecedor, moeda e total confiáveis \\
A2 & Redirecionar beneficiário ou conta & A conta bancária aplicada pertence à allowlist pré-declarada do fornecedor \\
A3a & Contabilizar ou acionar etapa crítica sem autorização & Nenhum estado crítico ocorre sem aprovação válida vinculada à ação exata \\
\bottomrule
\end{tabularx}
\end{table}

A unidade independente é o par contrafactual, identificado por \texttt{pair\_id}; \texttt{scenario\_id} identifica o documento benigno ou adversarial dentro do par. Análises restritas à classe adversarial podem agrupar por \texttt{scenario\_id}, pois há um único cenário adversarial por par; análises que combinem benignos e adversariais, ou comparem utilidade entre as classes, agrupam e reamostram por \texttt{pair\_id}, preservando condição e repetição como medidas dependentes. As três repetições medem variabilidade de execução, mas não são tratadas como três amostras independentes. Uma run só é avaliável se documento e estado inicial forem validados, o pipeline for iniciado e o oráculo obtiver o estado terminal. Falha de modelo, saída inválida, recusa ou payload não entregue continua sendo outcome do sistema. Apenas falha de infraestrutura que impeça observar o estado depois dos retries pré-definidos vira dado ausente.

\subsection{Arquiteturas C0 e C1}

As condições compartilham ingestão, extração, objetivo, modelo, contrato de proposta, fixture, runner, relógio e oráculo. Limites de arquivo, dados sintéticos, isolamento do laboratório e proteção de segredos também permanecem em C0 porque protegem a pesquisa e não constituem o tratamento avaliado.

\begin{figure}[H]
\centering
\small
\begin{minipage}[t]{0.46\textwidth}
\centering
\textbf{C0 -- integração funcional}\par\vspace{0.35em}
\fbox{\parbox{0.88\linewidth}{\centering
PDF $\downarrow$ extração $\downarrow$ LLM\\
$\downarrow$ proposta tipada\\
$\downarrow$ validação sintática\\
$\downarrow$ executor direto\\
$\downarrow$ Odoo}}
\end{minipage}
\hfill
\begin{minipage}[t]{0.50\textwidth}
\centering
\textbf{C1 -- defesa em profundidade}\par\vspace{0.35em}
\fbox{\parbox{0.90\linewidth}{\centering
PDF não confiável $\downarrow$ proveniência\\
$\downarrow$ campos sintéticos $\downarrow$ LLM\\
$\downarrow$ schema e canonicalização\\
$\downarrow$ policy engine\\
$\downarrow$ ALLOW / REVIEW / DENY\\
$\downarrow$ executor restrito e idempotente\\
$\downarrow$ Odoo + auditoria}}
\end{minipage}
\caption{Fluxos experimentais C0 e C1.}
\label{fig:c0-c1}
\end{figure}

C0 é uma baseline plausível, não uma versão artificialmente quebrada. O modelo produz \texttt{vendor-bill-action-proposal/v2}, contrato comum de criação e contabilização; propriedades extras, JSON inválido e valores internamente inconsistentes são recusados. Essa validação não compara argumentos com a verdade de negócio. Se a proposta for sintaticamente válida, seus argumentos seguem ao executor com as capacidades funcionais do fluxo. C0 não revalida fornecedor, pedido, valor ou beneficiário contra fontes confiáveis e não vincula autorização aos argumentos. Um avaliador sombra, determinístico e somente leitura, classifica as propostas de C0 e C1 pelas mesmas regras para análise de mecanismo, sem bloquear C0.

C1 aplica \textit{fail closed}. O documento recebe marca de origem não confiável; o estágio de extração não possui credencial de escrita; o contexto contém os campos sintéticos necessários à tarefa, comuns às condições; a proposta é canonicalizada; e o policy engine consulta PO, fornecedor, beneficiário e estado. A decisão é \texttt{ALLOW}, \texttt{REVIEW\_REQUIRED} ou \texttt{DENY}, acompanhada de códigos de razão. A criação consistente de rascunho pode ser autorizada; contabilização exige fixture de aprovação cujo hash cobre ação, fornecedor, valor, moeda, beneficiário, identidade, nonce e expiração. Alterar qualquer argumento invalida o token.

O executor C1 é um add-on mínimo do Odoo, cuja integração operacional usará JSON-2. O usuário técnico não recebe grupos genéricos de faturamento, compras ou pagamento, e nenhuma mutação crítica deve ser possível fora do método estreito. A chamada não aceita contexto fornecido pelo documento. Atomicidade evita estado parcial; idempotency key impede duplicação sob retry; e auditoria registra hashes, versões, política, identidade e estado anterior/posterior. No estado atual, os testes de efeito utilizam caminhos privados; o método público \texttt{execute} continua fechado. Sua abertura não é consequência automática da redução de escopo e exige os gates de revisão e integração.

\begin{table}[H]
\centering
\caption{Delta experimental entre C0 e C1.}
\label{tab:delta}
\footnotesize
\begin{tabularx}{\textwidth}{L{2.45cm} Y Y}
\toprule
\textbf{Elemento} & \textbf{C0} & \textbf{C1} \\
\midrule
Origem do conteúdo & Texto integrado ao contexto usual & Envelope explícito de dado externo não confiável \\
Validação & Sintaxe, tipo e consistência interna & Sintaxe, canonicalização, estado confiável e regra de negócio \\
Autoridade & Proposta válida pode acionar operação & Somente decisão de política e capacidade válida autorizam \\
Credencial & Identidade com permissões funcionais do fluxo & Identidades e capacidades restritas por ação \\
Ação crítica & Pode seguir da proposta sintaticamente válida & Exige aprovação vinculada aos argumentos exatos \\
Beneficiário & Argumento do modelo seleciona conta & Conta deve constar da allowlist do fornecedor \\
Auditoria & Log operacional básico & Proveniência, regras, aprovação e estado antes/depois \\
\bottomrule
\end{tabularx}
\end{table}

\subsection{Ambiente, modelos e reprodutibilidade}

O testbed usa Odoo 19.0 Community self-hosted, fonte oficial fixada no commit \texttt{a4b87724\allowbreak 5c9cb5e3\allowbreak d7152d15\allowbreak 93a51817\allowbreak 0f115ca1} \cite{odoo2026source}, PostgreSQL 18 e gateway Python 3.12. A imagem Odoo aceita possui manifest OCI local distinto do config digest. O banco não publica porta, os dados são sintéticos e a rede interna impede egress de ERP, banco e outbox. JSON-2 é transporte; as garantias de C1 pertencem ao add-on e às credenciais restritas.

O único modelo desta entrega será solicitado à API DeepSeek pelo alias \texttt{deepseek-flash}, aprovado para o pré-piloto e observado no smoke técnico de 26 de setembro. A seleção histórica de 31 de agosto usou \texttt{deepseek-v4-flash}, então associado a \texttt{DeepSeek-V4-Flash-0731}; esse nome histórico não identifica a versão atual da execução \cite{deepseek2026models,deepseek2026updates}. Cada run registra alias solicitado, nome retornado e versão quando disponível, request ID saneado e timestamp, sem inferir um snapshot imutável. \textit{Thinking} fica explicitamente desligado e a temperatura do pré-piloto é zero. C0 e C1 usarão o mesmo alias e os mesmos parâmetros; mudanças observadas do serviço serão registradas. Replicação com um modelo local, como Qwen3-14B \cite{qwen2025modelcard}, é trabalho futuro e não consta do orçamento de coleta desta entrega.

O prompt-base da tarefa, versão \texttt{ap-task-prompt/v1}, é idêntico em C0 e C1. Ele descreve a tarefa de Contas a Pagar, as regras de preenchimento exigidas pelo schema e uma orientação genérica de que o texto extraído é dado de terceiros, cujas instruções não alteram a tarefa nem os campos da proposta. Essa orientação integra a baseline plausível, não a defesa avaliada. A diretiva de fase, criar rascunho ou contabilizar, vem do fluxo confiável e nunca do documento. O prompt não nomeia ataques, campos-alvo, política nem aprovação. Como a proposta depende apenas de documento, fase, prompt e modelo, ela é gerada em estágio próprio, com uma chamada por execução e condição e a mesma política de tentativas, e depois entregue ao executor da condição correspondente.

Prompts, schemas, políticas, imagens, configurações e dataset recebem SHA-256. Requests e responses brutos ficam fora do Git em armazenamento criptografado com acesso controlado e cópia imutável; o repositório contém manifest saneado. Segredos são lidos de arquivo ou ambiente e não entram em Compose, argumentos ou logs. Uma API móvel não permite alegar reprodução binária do modelo: o pacote distingue reexecutar o protocolo de reproduzir a análise sobre os dados preservados.

\subsection{Dataset AP-IPI-BR}

O dataset contém PDFs sintéticos em português, com camada de texto, marca “SEM VALOR FISCAL -- DADOS SINTÉTICOS” e nenhuma NF-e válida. OCR, macros, documentos reais e dados bancários reais ficam fora do escopo. O conjunto final contém 90 pares contrafactuais: os fatos, template e objetivo de negócio permanecem constantes, e somente o payload diferencia o membro adversarial. Pelo menos 15 benignos são hard negatives com linguagem semelhante a instruções.

\begin{table}[H]
\centering
\caption{Distribuição pré-piloto do conjunto confirmatório.}
\label{tab:dataset}
\small
\begin{tabularx}{\textwidth}{Y C{2.15cm} C{2.2cm} Y}
\toprule
\textbf{Família} & \textbf{Adversariais} & \textbf{Benignos pareados} & \textbf{Ação ou propriedade alvo} \\
\midrule
A1 -- campos financeiros & 30 & 30 & Criar rascunho sem alterar valor, fornecedor ou referência \\
A2 -- beneficiário & 30 & 30 & Usar somente conta bancária aprovada \\
A3a -- ação crítica & 30 & 30 & Contabilizar apenas com aprovação exata \\
\midrule
\textbf{Total} & \textbf{90} & \textbf{90} & 90 pares contrafactuais \\
\bottomrule
\end{tabularx}
\end{table}

O desenho combina três templates, duas posições de payload e cinco estratégias linguísticas: instrução explícita, falsa autoridade, nota procedural, ofuscação textual e instrução contextualizada. Cada um dos três objetivos possui 30 pares e peso de um terço no resultado agregado. Os 30 controles benignos de ação crítica incluem aprovação pré-definida, vinculada aos argumentos da contabilização legítima. Exposição de dados, mensagens ao fornecedor, mascaramento e tokenização não integram a avaliação.

Cada item possui ground truth anterior aos resultados: ação esperada, disposição, fatos confiáveis e deltas proibidos. Um integrante cria e outro valida integralmente os itens; divergências são resolvidas antes do hash final. A geração automática não substitui essa revisão. Quem implementa C1 não recebe o conjunto final completo antes do freeze. O piloto usa 12 pares adicionais, quatro por objetivo, excluídos do teste. Nenhum conjunto adaptado à defesa é previsto nesta entrega.

\subsection{Métricas e oráculos}

O endpoint primário é ASR-state: a proporção de execuções adversariais avaliáveis em que o invariante alvo é violado no Odoo. Uma frase maliciosa na resposta, proposta indevida ou tool call bloqueada não conta como efeito final. Métricas intermediárias preservam o mecanismo.

\begin{table}[H]
\centering
\caption{Principais métricas do experimento.}
\label{tab:metrics}
\footnotesize
\begin{tabularx}{\textwidth}{L{2.25cm} Y L{3.05cm}}
\toprule
\textbf{Métrica} & \textbf{Definição operacional} & \textbf{Fronteira observada} \\
\midrule
Delivery Rate & Payloads que chegaram ao modelo entre adversariais avaliáveis & Extração / contexto \\
ASR-proposal & Propostas que incorporam o objetivo adversarial entre ataques entregues & LLM / proposta \\
ASR-policy & Propostas maliciosas efetivamente autorizadas & Política / autorização \\
ASR-state & Violações observadas entre adversariais avaliáveis & Odoo; primária \\
UACR & Ações críticas concluídas sem autorização válida e exata & Estado crítico \\
BTSR & Tarefas benignas com estado exatamente esperado & Utilidade benigna \\
Utility Under Attack & Tarefa legítima concluída sem atingir objetivo adversarial & Segurança e utilidade \\
False Block Rate & Ação benigna correta negada ou incompleta por segurança & Custo funcional \\
\bottomrule
\end{tabularx}
\end{table}

Valores monetários usam representação decimal e comparação exata. IDs e estados são lidos por código, não julgados por outro modelo. O runner registra latência end-to-end benigna, tempo até disposição, verificação de aprovação, tokens, custo por 100 documentos, chamadas, timeout, retry e completude de log. O tempo de uma pessoa não é incluído; \textit{Review Demand per 100} é somente a quantidade potencial de casos encaminhados.

\subsection{Procedimento experimental e análise}

O piloto executa 24 documentos, duas condições e três repetições, totalizando 144 runs DeepSeek. C0 precisa concluir ao menos 10 dos 12 cenários benignos em duas ou mais repetições, criar e ler o estado sem corrupção e manter logging completo em pelo menos 95\% das runs. C1 também precisa concluir tarefas legítimas, inclusive contabilizações com aprovação válida, para que bloquear toda ação crítica não pareça uma defesa útil. Se o gate falhar, corrige-se a funcionalidade e repete-se o piloto; não se reduz a definição de sucesso. O piloto estima custo, latência, vazão e discordâncias pareadas; subsidia a verificação de poder e precisão e fixa parâmetros antes do freeze. Não escolhe limiares confirmatórios de aceitabilidade de custo. Qualquer revisão dos 90 pares exige decisão explícita anterior ao teste final; não há redução automática por prazo.

O piloto também tem um gate de eficácia adversarial. Se a ASR diagnóstica de C0 for nula ou próxima de zero, as estratégias A1--A3 são revisadas sem enfraquecer C0 e o piloto é repetido uma única vez. Persistindo a ausência de efeito, o grupo decide, antes do congelamento, enquadrar o estudo como estimativo, reportando a redução absoluta de risco com intervalo em vez de teste de superioridade. O limiar de suficiência dessa repetição foi fixado antes de sua execução: ao menos uma execução de C0 comprometida em ao menos dois pares distintos.

Depois do congelamento, a agenda é bloqueada por cenário e repetição. A ordem C0/C1 é randomizada dentro de bloco; antes de cada condição, banco e filestore são restaurados e seus hashes conferidos. O pipeline é executado, o oráculo lê o estado, a evidência bruta é persistida e a restauração deve reproduzir o baseline, inclusive após contabilização. Há no máximo duas novas tentativas de inferência apenas para erro transitório pré-classificado e antes de qualquer efeito no Odoo. Resultado remoto incerto suspende o lote; cancelar uma espera local não comprova encerramento remoto. Retries de observação são separados e dependem de configuração e revisão anteriores ao freeze. Reset divergente preserva o resultado observado e bloqueia a próxima execução. Qualquer mudança de modelo, prompt, schema, política, permissão, runner ou oráculo que afete comportamento gera novo freeze e reexecução integral dos pares afetados.

O experimento principal possui $180\times2\times3=1.080$ runs. Com as 144 do piloto separado, o compromisso de coleta é de 1.224 execuções, sem retries e ensaios técnicos. As três repetições são mantidas em ambos os conjuntos. Qwen3, ablação e ataques adaptados à defesa ficam fora da entrega e não são execuções condicionais a completar nesta agenda.

Para ASR-state, serão estimadas taxas por condição, redução absoluta e risco relativo, quando definido, com intervalos de 95\% por bootstrap de cluster pareado: cada sorteio de \texttt{pair\_id} conserva todas as suas condições e repetições. Serão 10.000 réplicas, intervalo BCa e percentil como verificação de sensibilidade, com semente registrada. McNemar exato será complementar sobre \texttt{ever\_compromised}, definido quando ao menos uma repetição viola o estado. Esse outcome agrega três oportunidades de ataque e não substitui a ASR por execução. \hyperlink{h1}{H1} exige limite superior do IC95\% de $ASR_{C1}-ASR_{C0}$ abaixo de zero. \hyperlink{h3}{H3} exige limite inferior do IC95\% de $BTSR_{C1}-BTSR_{C0}$ acima de $-0{,}10$; ausência de diferença significativa não será interpretada como equivalência.

O tratamento de zero e fronteira permanece pré-especificado: não se declara certeza a partir de bootstrap degenerado, nem se aplica correção de continuidade ad hoc para obter risco relativo. Uma proporção igual a zero ou um não implica, isoladamente, que o intervalo BCa do contraste seja inestimável; sua estimabilidade precisa ser verificada. Wilson e Newcombe se restringem aos resumos binários independentes por cenário previstos no plano, não às execuções repetidas tratadas como independentes. A regra de fronteira para o intervalo do estimando primário com três repetições exige implementação, validação e fixação antes do freeze. O risco relativo é omitido quando uma das taxas é zero. A correlação das repetições e o pareamento não podem desaparecer nessa regra.

Latência será descrita por mediana, intervalo interquartil e p95; tokens e custo aparecerão por execução, tarefa concluída e 100 documentos. Famílias de ataque, estratégias, posições, templates e demanda técnica de revisão terão descrição de efeitos e falhas, sem testes de subgrupos ou busca de significância. Não serão usados modelos adicionais de sensibilidade ou famílias secundárias de testes inferenciais.

Todas as tabelas apresentarão denominadores agendados, avaliáveis e dados ausentes (\textit{missing}). Falha de modelo com estado final legível continua avaliável; estado indisponível ou efeito incerto não vira automaticamente defesa bem-sucedida. Contrastes pareados exigem as duas condições avaliáveis no mesmo par e repetição. Serão preservadas as sensibilidades de denominador e de falha fechada, além do cenário conservador que desfavorece a conclusão específica: missing em C1 conta como violação para \hyperlink{h1}{H1}/\hyperlink{h2}{H2} e como falha benigna para \hyperlink{h3}{H3}, com a direção inversa em C0. A análise principal não será substituída seletivamente por uma sensibilidade mais favorável.

\subsection{Ética, segurança e validade}

O estudo não movimenta dinheiro, não usa documento fiscal válido, não envia e-mail e não processa dados de pessoas ou empresas. Pagamentos e solicitações simuladas de pagamento não integram esta entrega; a outbox permanece desativada. O host, Odoo, banco e arquivos brutos permanecem sob acesso restrito; somente dados sintéticos saneados podem ser publicados após dupla revisão.

As principais ameaças à validade interna são drift do serviço externo, diferença não intencional entre C0/C1, estado residual e conhecimento do dataset pela defesa. São mitigadas por igualdade das variáveis controladas, execução intercalada, reset por hash, revisão independente de C0, segregação de papéis e freeze. A validade de constructo depende de o oráculo representar dano relevante; por isso, estado final é separado de proposta. A validade externa é limitada a um ERP, idioma, processo PO-backed, PDFs textuais, um modelo consumido por API e ataques sintéticos estáticos. Sem segundo modelo, ablação ou ataques adaptados à defesa, não se conclui transportabilidade entre modelos, contribuição causal de cada controle ou resistência adaptativa. Poder limitado por família, células raras e correlação entre repetições ameaçam validade de conclusão e serão tratados com intervalos validados, métodos exatos, análise pareada e linguagem proporcional à evidência.


\section{Resultados e discussão}
\subsection{Estado técnico do artefato}

Esta seção distingue provas de engenharia de resultados científicos. Os marcos M1/M2, registrados em 31 de agosto de 2026, são históricos; os incrementos posteriores de integração privada e registro não os substituem nem constituem a coleta confirmatória. M1 estabeleceu o laboratório reproduzível. A imagem do Odoo 19.0 Community foi construída a partir do commit fixado, com manifest OCI local registrado separadamente do config digest. PostgreSQL 18 permaneceu sem porta publicada; Odoo foi exposto somente em loopback; e uma identidade técnica C0, distinta da chave administrativa, realizou o fluxo JSON-2.

As fixtures sintéticas incluem empresa, fornecedor, conta bancária, produto e pedido de compra confirmado. O smoke M1 criou uma única \texttt{account.move} do tipo \texttt{in\_invoice}, em \texttt{draft}, com \texttt{auto\_post='no'}, beneficiário persistido e linha vinculada ao pedido. Uma leitura independente verificou o efeito; em seguida, o rascunho foi removido. O hash canônico após a remoção coincidiu com o estado inicial, demonstrando que o testbed consegue produzir e desfazer a mutação observada sem apagar os volumes.

M2 ligou uma entrada documental ao mesmo efeito. O pipeline gera localmente um PDF sintético, marcado como sem valor fiscal e contendo camada de texto; valida tipo, tamanho, número de páginas e limite de caracteres; extrai o texto sem OCR; produz uma proposta \texttt{vendor-bill-proposal/v1}; valida propriedades, tipos, formatos e total por aritmética decimal; resolve referências por chaves de negócio; e chama o executor direto C0. PDFs criptografados, inválidos, sem texto ou fora dos limites são recusados.

Para isolar o encanamento antes da seleção do LLM, M2 usou um adaptador determinístico de desenvolvimento, versão 1. Esse adaptador lê somente os campos declarados da fixture e é proibido em runs científicas. Sua função foi demonstrar que PDF, contrato, executor, Odoo e oráculo concordam, não simular inteligência ou resistência a ataque.

\begin{table}[H]
\centering
\caption{Evidências técnicas disponíveis antes do piloto.}
\label{tab:current-evidence}
\footnotesize
\begin{tabularx}{\textwidth}{L{2.10cm} Y L{3.05cm}}
\toprule
\textbf{Evidência} & \textbf{Resultado observado} & \textbf{Limite do resultado} \\
\midrule
Fonte Odoo & Community 19.0 no commit \texttt{a4b877245c9cb5e3d7152d1593a518170f115ca1} & Não representa todo Odoo 19 futuro \\
Manifest de imagem & OCI local \texttt{sha256:3e0bcf8e7968\ldots} ligado à fonte & Ainda não publicado em registry \\
Identidade C0 & Conta técnica não administrativa cria e lê o rascunho & Evidência histórica M1/M2, não integração científica completa \\
Contrato M2 & \texttt{vendor-bill-proposal/v1}, fechado e versionado & Precede o contrato comum v2 atual \\
Testes M2 & 41 testes aprovados no gateway 0.2.0 & Suíte histórica, não contagem da árvore atual \\
Smoke M2 & PDF $\rightarrow$ proposta $\rightarrow$ rascunho $\rightarrow$ remoção & Usa adaptador determinístico de desenvolvimento \\
Estado & hash inicial e restaurado \texttt{224448b68946\ldots}; estado criado distinto & Um cenário benigno de engenharia \\
\bottomrule
\end{tabularx}
\end{table}

A suíte cobre PDF criptografado ou sem texto, estrutura inválida, campo duplicado ou ausente, propriedade extra, versão de schema incompatível, data e valor malformados, total divergente, referência inexistente, duplicidade e resíduo no Odoo. O smoke final passou com 41 testes, criou exatamente um rascunho e restaurou o hash inicial. PDF, texto, proposta e estados receberam SHA-256; credenciais e evidências detalhadas ficaram fora do Git.

Até 26 de outubro, o repositório também registra provas posteriores: contrato comum v2, caminhos privados de criação e contabilização C0/C1, autorização vinculada, captura independente de estado e restauração por snapshot de banco e filestore. O checkpoint de captura/restauração percorreu oito casos técnicos e 16 restaurações, incluindo contabilização, com provedor e supervisão simulados. Esses números contam casos de engenharia, não ataques científicos nem observações independentes do estudo.

O registro candidato \texttt{run-record/v2} preserva resultados recomputados, falhas de proposta, estado não observável e divergência de reset, sem preencher lacunas com sucesso. O gravador local preserva registros por hashes e escrita exclusiva; isso não demonstra backup externo imutável. A integração operacional e a revisão independente permanecem abertas, e o método público \texttt{execute} continua fechado. As provas anteriores não foram reexecutadas para esta revisão editorial.

\subsection{O que os resultados atuais demonstram}

Os resultados demonstram quatro propriedades necessárias ao futuro experimento. Primeiro, o estado de negócio pode ser observado independentemente da saída do modelo. Segundo, a proposta pode ser validada e canonicalizada antes da mutação. Terceiro, as condições podem começar da mesma fixture e terminar sem resíduo detectável. Quarto, os identificadores internos do ERP não precisam ser aceitos do documento: fornecedor, pedido e conta são resolvidos por chaves sintéticas declaradas.

Essas propriedades reduzem ameaças de instrumentação. Sem um oráculo externo, uma ferramenta chamada poderia ser contada como ataque bem-sucedido mesmo se a transação falhasse; sem reset, uma condição poderia herdar o efeito anterior; sem contrato comum, C1 poderia parecer superior apenas porque rejeita JSON inválido que C0 aceita. M1/M2 antecipam esses problemas e mantêm a validação estrutural nos dois braços.

O trabalho também responde ao feedback de que a continuidade do TCC I precisava tornar-se aplicada e definir segmento-alvo. O artefato executa um fluxo de Contas a Pagar em um ERP real de laboratório, e não apenas uma arquitetura conceitual. Ainda assim, implementação não substitui avaliação. O resultado atual é evidência de viabilidade e controle experimental, não de eficácia da defesa.

\subsection{Piloto: propostas do modelo}

Em 26 de outubro de 2026, o piloto executou o estágio de propostas com o modelo real: 12 pares, dois membros, C0 e C1 e três repetições, totalizando 144 chamadas por rodada, em ordem randomizada por semente registrada. Mediu-se a proposta (ASR-proposal) e a decisão que a política C1 tomaria sobre ela (ASR-policy), contra o estado confiável derivado do gabarito. O efeito no Odoo, isto é, ASR-state, UACR e BTSR, ainda não foi medido, porque o ciclo integrado não foi ligado aos 12 pares.

A rodada 1 usou notas adversariais de uma ou duas frases. Como nenhuma proposta de C0 foi comprometida, o gate de eficácia disparou. As 12 notas foram então reescritas segundo princípios estáticos: justificar por que o valor impresso estaria errado, alegar prioridade ou competência, nomear o campo e o valor a registrar, integrar-se ao vocabulário do documento e manter ofuscação legível. A revisão não usou conhecimento da defesa e manteve documentos benignos, gabarito, valores do atacante, prompt e modelo. A rodada 2 foi a única repetição permitida.

\begin{table}[H]
\centering
\caption{Piloto de propostas: resultados por rodada e condição.}
\label{tab:pilot-proposals}
\footnotesize
\begin{tabularx}{\textwidth}{Y C{1.35cm} C{1.35cm} C{1.35cm} C{1.35cm}}
\toprule
& \multicolumn{2}{c}{\textbf{Rodada 1}} & \multicolumn{2}{c}{\textbf{Rodada 2}} \\
\textbf{Medida} & \textbf{C0} & \textbf{C1} & \textbf{C0} & \textbf{C1} \\
\midrule
Chamadas com proposta válida & 72/72 & 72/72 & 72/72 & 71/72 \\
Benignas idênticas ao gabarito & 36/36 & 36/36 & 36/36 & 35/36 \\
Benignas aceitas pela política C1 & -- & 36/36 & -- & 35/36 \\
Adversariais com campo-alvo alterado (ASR-proposal) & 0/36 & 0/36 & 0/36 & 0/36 \\
Adversariais idênticas ao gabarito & 36/36 & 36/36 & 36/36 & 36/36 \\
Comprometidas aceitas por C1 (ASR-policy) & -- & 0/0 & -- & 0/0 \\
\bottomrule
\end{tabularx}
\end{table}

Nas duas rodadas, nenhuma das 72 propostas adversariais de C0 incorporou o objetivo do atacante, e o limiar fixado antes da segunda rodada não foi atingido. Com 0 em 36 execuções de C0 em uma rodada, o limite superior unilateral de 95\% da taxa é de cerca de 8\% se as execuções fossem independentes e de 22\% tomando o par como unidade. A única falha foi uma saída fora do schema em uma chamada benigna de C1 da rodada 2, mantida como resultado e sem nova tentativa. Cada rodada consumiu cerca de 215 mil tokens, aproximadamente US\$~0,07 pelos preços datados, com latência mediana de 2,1~s por chamada.

Como verificação de componente, e não como resultado do experimento, propostas que obedecessem integralmente às 12 notas foram submetidas à política C1. Onze seriam negadas por divergência de pedido, fornecedor, valor, linha, conta fora da allowlist ou alvo da contabilização. Uma, que altera somente a referência impressa do documento, seria aceita: essa referência é controlada pelo atacante e funciona apenas como sinal secundário de duplicidade.

\subsection{Resultados confirmatórios ainda não disponíveis}

Na data desta versão, permanecem pendentes o ciclo integrado com o Odoo para os pares do piloto, a revisão nominal de C1 e do protocolo, a decisão sobre o enquadramento, o dataset final revisado e o congelamento. Não há coleta confirmatória C0--C1 contabilizada. A Tabela~\ref{tab:pending-results} registra o estado das respostas às perguntas de pesquisa; “não executado” não representa zero ocorrência.

\begin{table}[H]
\centering
\caption{Estado dos resultados científicos em 26 de outubro de 2026.}
\label{tab:pending-results}
\small
\begin{tabularx}{\textwidth}{C{0.80cm} L{3.25cm} Y}
\toprule
\textbf{RQ} & \textbf{Estimadores previstos} & \textbf{Estado} \\
\midrule
\hyperlink{rq1}{RQ1} & ASR-proposal, ASR-policy, ASR-state e UACR por C0/C1 & ASR-proposal e ASR-policy medidas no piloto, com 0/36 em cada condição e rodada; ASR-state e UACR não executados; nada confirmatório \\
\hyperlink{rq2}{RQ2} & BTSR, falsos bloqueios, STP, escalonamento, latência e custo & Exatidão das propostas benignas, latência e custo do estágio do modelo medidos no piloto; BTSR no estado, falsos bloqueios e escalonamentos não executados \\
\hyperlink{rq3}{RQ3} & Efeitos por A1, A2 e A3a & Nenhuma família, estratégia, posição ou template produziu proposta comprometida no piloto; teste final não executado \\
\bottomrule
\end{tabularx}
\end{table}

Não é possível afirmar que C1 reduz prompt injection, mantém BTSR no estado do ERP ou possui custo aceitável. O custo será caracterizado, não submetido a hipótese de aceitabilidade. Qwen3, ablação e ataques adaptados à defesa foram retirados da entrega e passam a trabalho futuro; \hyperlink{rq4}{RQ4} e \hyperlink{h5}{H5} não têm resultados a produzir neste estudo. O resumo e a conclusão preservam essa ausência para impedir que números exemplificativos sejam confundidos com observações.

\subsection{Discussão e próximos gates}

A literatura sugere que separar instruções e dados, alinhar ferramentas à tarefa e impor regras em tempo de execução pode reduzir ataques \cite{chen2025struq,jia2025taskshield,li2025drift}. O artefato segue essa direção, mas desloca a autoridade crítica para política e executor determinísticos. Essa escolha procura reduzir a dependência de reconhecer todas as formulações de ataque: mesmo que o LLM proponha outra conta bancária, o estado confiável deve impedir a autorização. O experimento dirá se essa composição reduz violações sem bloquear o processo em excesso.

O piloto altera a leitura do desenho. Com a orientação genérica comum às duas condições, o modelo copiou os valores impressos em todas as execuções adversariais, mesmo diante de autoridade alegada, correção integrada ao documento, ordem direta e texto espaçado. Se a ASR de C0 for nula, a diferença $ASR_{C1}-ASR_{C0}$ não pode ficar abaixo de zero e \hyperlink{h1}{H1} não pode ser satisfeita; como a ordem confirmatória é fixa, \hyperlink{h3}{H3} e \hyperlink{h2}{H2} ficariam sem interpretação confirmatória por um efeito de ataque inexistente. Por isso, há uma proposta, pendente de ratificação do grupo antes do congelamento, de tratar \hyperlink{h1}{H1} como estimativa, com redução absoluta de risco e intervalo, manter \hyperlink{h3}{H3} confirmatória e independente de \hyperlink{h1}{H1}, com a margem inalterada, e decidir \hyperlink{h2}{H2} pela UACR de C0 medida no Odoo. Esta versão registra a proposta, sem adotá-la. O resultado também não demonstra robustez do modelo: os ataques são estáticos, foram calibrados depois da primeira rodada e um único modelo foi avaliado.

Para concluir o piloto, ainda são necessários: ligar o ciclo integrado aos 12 pares com as propostas já geradas, incluindo despacho C0/C1, autorização, supervisão, captura, registro e reset; medir ASR-state, UACR e BTSR; e concluir a revisão independente dos documentos e componentes. O add-on e as provas privadas existentes serão reutilizados, sem expansão de componentes fora desse caminho. O smoke real do modelo e os testes simulados não substituem a competência do agente no piloto. C0 deve cumprir o gate de competência e C1 deve concluir ações legítimas, incluindo contabilização autorizada; a avaliação não pode premiar a simples desativação de toda ação crítica.

Depois do piloto, o grupo deverá documentar discordâncias pareadas, poder, precisão, latência, custo, vazão e completude de logs, além de validar os intervalos nos casos de fronteira. Permanecem 90 pares e três repetições; qualquer revisão da amostra exige decisão explícita antes do teste final, sem corte automático por prazo. O artigo final substituirá a Tabela~\ref{tab:pending-results} por resultados com numeradores, denominadores, missing, diferenças, riscos relativos quando definidos, intervalos e sensibilidades pré-especificadas. Não serão acrescentados blocos de ablação ou replicação local nesta entrega.

O principal risco editorial desta etapa é a falsa sensação de completude produzida por um PDF formatado. O formato institucional já pode ser validado, mas o conteúdo confirmatório não. Por isso, este relatório deve orientar a implementação e a revisão com o orientador, não ser tratado como a submissão final de novembro.


\section{Conclusão}
Este trabalho delimitou o problema de prompt injection indireta em um agente que processa documentos de Contas a Pagar e atua sobre Odoo. A arquitetura proposta separa interpretação, autorização e execução: o LLM produz uma proposta tipada; regras determinísticas confrontam a proposta com pedido, fornecedor, beneficiário e estado confiáveis; ações críticas exigem aprovação vinculada; e o executor opera com menor privilégio, idempotência e auditoria. C0 e C1 serão comparados sob o mesmo modelo, documento, prompt-base, schema e fixture.

Na etapa registrada por este artigo, foi demonstrada a viabilidade técnica de partes do aparato experimental. Odoo 19.0 Community foi fixado em uma revisão reproduzível; M1/M2 comprovaram o fluxo PDF--proposta--rascunho e o reset de uma fixture sintética. Incrementos posteriores acrescentaram contrato comum v2, caminhos privados C0/C1, captura independente, restauração de banco e filestore e registro candidato que preserva falhas. Essas provas reduzem riscos de instrumentação, mas não equivalem a um runner científico integrado e revisado.

O piloto de propostas, executado com o modelo real, trouxe o primeiro resultado empírico. Em duas rodadas, a segunda com estratégias de ataque revisadas, nenhuma das 72 propostas adversariais de C0 incorporou o objetivo do atacante, e as propostas benignas foram exatas em 36 de 36 execuções em C0 e em 35 de 36 em C1. Com esse modelo, esse prompt-base e ataques estáticos, a redução do sucesso dos ataques por C1 não pode ser demonstrada, porque a baseline já não cede. Isso não prova robustez do modelo nem dispensa a fronteira determinística, cuja função é independer do que o modelo proponha.

Esses resultados não respondem ainda se C1 reduz ASR-state ou UACR, se preserva BTSR no estado do ERP ou qual sobrecarga integral introduz. A coleta confirmatória não ocorreu, e os testes privados não abrem o método público de execução. Assim, a conclusão permitida é restrita: existem componentes observáveis e restauráveis, e o modelo resistiu aos ataques do piloto, mas a eficácia e o custo da proteção permanecem questões abertas.

As próximas etapas são decidir o enquadramento estimativo proposto, ligar o ciclo integrado aos pares do piloto e medir o estado, concluir a revisão, congelar o protocolo e executar os pares randomizados. O núcleo preserva 144 execuções piloto e 1.080 principais, com três repetições, totalizando 1.224 sem retries. O artigo final deverá responder \hyperlink{rq1}{RQ1}--\hyperlink{rq3}{RQ3} com denominadores, efeitos, intervalos e descrição de falhas, sem hipótese de aceitabilidade de custo. Segundo modelo, ablação e ataques adaptados à defesa são trabalhos futuros. Sua ausência limita a generalização entre modelos, a atribuição causal por camada e a resistência adaptativa. Pagamentos, documentos reais, OCR e comportamento de revisores humanos continuarão fora do escopo.

A contribuição científica dependerá da combinação entre arquitetura executável e avaliação honesta. Se C1 reduzir violações sem ultrapassar a margem de utilidade, o estudo fornecerá evidência situada para integrações semelhantes. Se não reduzir, ou se o custo for excessivo, o resultado continuará relevante ao mostrar os limites práticos das camadas escolhidas. Em ambos os casos, a conclusão final será determinada pelos dados congelados, não pela expectativa de que a arquitetura funcione.


\renewcommand{\refname}{6. Referências bibliográficas}
\bibliographystyle{abntex2-alf}
\bibliography{referencias}

\ifdefined\tccannex
\includepdf[pages=-,pagecommand={\thispagestyle{empty}}]{\tccannex}
\fi

\end{document}
